\chapter{K-stability for K\"ahler models}
\index{K-stability!Kahler model@K\"ahler model}\index{cscK metric}\index{Mabuchi functional!non-Archimedean}
\label{chap:mpwn-k-stability}

The preceding chapters provide two ways to test positivity at infinity:
intersection numbers of analytic degenerations, and coercivity of the
non-Archimedean Mabuchi functional on all finite-energy potentials.  The
transcendental Calabi--Yau theorem relates these tests through big models,
continuous envelopes, and divisorial measures.  Positive-product formulas
require a normalization of the model divisor.  We keep that hypothesis
visible and distinguish the normalized model inequality from the literal
test over every big total-space class.

Throughout, $X$ is a connected compact K\"ahler manifold of dimension
$n\geq1$, and $\alpha$ is a K\"ahler class.  In every use of a
positive-product identity with a potential, including the coefficient
comparison below, the model divisor satisfies $f_D\geq1$.
Put
% RETHLAS-VERIFIED: mpwn-factor-n-correction; receipt 20260814T175959Z_484442fcc393
\[
                    V:=\alpha^n,\qquad
                    s:=-nV^{-1}\alpha^{n-1}\cdot K_X.
\tag{62.1}\label{eq:mpwn-s-constant}
\]
Positive products are written \(\langle\ \rangle\).  They agree with
ordinary intersections for K\"ahler classes, but not for arbitrary big
classes.  The factor \(n\) in \(s\) is the scalar-curvature normalization
compatible with the Donaldson--Futaki coefficient below.  The formula for
\(s\) printed in \cite[Section~2.4.2]{MPWN25} omits this factor, whereas
\cite[Definition~1.1.1 and Theorem~9.1.1]{MPWN25}, together with the
normalization \(E_A=(n+1)^{-1}V^{-1}(A,\phi)^{n+1}\), force the displayed
value.  Here is the exact coefficient comparison.  On a big model, put
\[
 B:=A+D,\qquad C:=\alpha^{n-1}\cdot K_X,\qquad
 Q:=\langle B^{n+1}\rangle
\]
and
\[
 R(B):=K^{\log}_{\cX/\mathbf P^1}\cdot\langle B^n\rangle,\qquad
 K^{\log}_{\cX/\mathbf P^1}
 :=K_{\cX/\mathbf P^1}-(\cX_0-\cX_{0,\mathrm{red}}).
\]
The model formulas for energy, twisted energy, and reduced entropy give
\[
 V M_\alpha^{\NA}(P_\alpha(f_D))
       =R(B)+\frac{s}{n+1}Q,
\tag{62.1a}\label{eq:mpwn-factor-n-potential-side}
\]
whereas the Donaldson--Futaki and model Mabuchi definitions give
\[
 M^{\NA}(\cX,B)
       =R(B)-\frac{nC}{(n+1)V}Q.
\tag{62.1b}\label{eq:mpwn-factor-n-model-side}
\]
Thus the twisted-energy and entropy terms contribute the common term
\(R(B)\), not another multiple of \(Q\).  To see that coefficient comparison
is legitimate, take the product model and
\(D=c\cX_0\) with \(c\geq1\).  Then \(f_D=c\geq1\), the class
\(A+c[\cX_0]\) is K\"ahler, and
\[
       Q=(A+c[\cX_0])^{n+1}=(n+1)cV\ne0.
\]
The model identity in Theorem~9.1.1 of \cite{MPWN25} therefore forces
\(s=-nC/V\), which is \eqref{eq:mpwn-s-constant}.

This source-text correction was independently verified by
\texttt{gpt-5.6-sol} at \texttt{xhigh} scrutiny.  The receipt has verdict
\texttt{correct}, with no gaps; its exact identifier, digest, and provenance
are recorded in the editorial ledger and \Cref{app:provenance}.

\section{Analytic test configurations and big models}

\begin{definition}[Smooth dominating analytic test configuration]
\label{def:mpwn-analytic-test-configuration}
A test configuration for \((X,\alpha)\) consists of:
\begin{enumerate}
\item a compact K\"ahler manifold \(\cX\) and a
      \(\mathbf C^*\)-equivariant proper bimeromorphic morphism
      \[
            \mu:\cX\longrightarrow X\times\mathbf P^1
      \tag{62.2}\label{eq:mpwn-model-map}
      \]
      that is a biholomorphism away from the zero fiber;
\item the lifted standard \(\mathbf C^*\)-action;
\item a class
      \[
                 \mathcal A_D:=A+[D]\in H^{1,1}_{\mathrm{BC}}(\cX,\mathbf R),
          \qquad A:=\mu^*\operatorname{pr}_1^*\alpha,
      \tag{62.3}\label{eq:mpwn-model-class}
      \]
      with \(D\) a real vertical divisor supported on \(\cX_0\).
\end{enumerate}
It is a K\"ahler test configuration if \(\mathcal A_D\) is K\"ahler and a big
test configuration if \(\mathcal A_D\) is big.
\end{definition}

This is \cite[Definition~1.1.1]{MPWN25}, in the smooth dominating form used
for intersection computations.  An equivariant resolution makes this
restriction cofinal.  A class is big when it contains a K\"ahler current;
it need not be nef.  Therefore a big test configuration does not, by
itself, define a model psh function.  Its correct potential is the envelope
\(P_\alpha(f_D)\).

\begin{example}[Deformation to the normal cone]
\label{ex:mpwn-dnc-stability}
For a smooth submanifold \(Z\subset X\), blow up
\(Z\times\{0\}\subset X\times\mathbf P^1\), with exceptional divisor \(E\).
The classes \(A-c[E]\) are relatively K\"ahler for sufficiently small
\(c>0\).  Equivalently, after adding a sufficiently large multiple of the
central fiber, they are K\"ahler on the total space.  At the boundary of
the relative K\"ahler range they are relatively nef, and beyond it they
may remain relatively big.  Global bigness can be obtained by adding
a base class; the effect of this addition on positive-product expressions
must be checked, as in \Cref{ex:mpwn-big-translation-failure}.
Testing only the small-$c$ interval misses the positive-intersection
data carried by the big range.
\end{example}

\section{Donaldson--Futaki, Mabuchi, and
\texorpdfstring{\(J\)}{J}}

For a big test configuration, define
\[
 \DF(\cX,\mathcal A_D)
 :=K_{\cX/\mathbf P^1}\cdot\langle\mathcal A_D^n\rangle
   -\frac{n\,\alpha^{n-1}\cdot K_X}{(n+1)\alpha^n}
       \langle\mathcal A_D^{n+1}\rangle,
\tag{62.4}\label{eq:mpwn-df}
\]
\[
 M^{\NA}(\cX,\mathcal A_D)
 :=\DF(\cX,\mathcal A_D)
   -(\cX_0-\cX_{0,\mathrm{red}})
          \cdot\langle\mathcal A_D^n\rangle,
\tag{62.5}\label{eq:mpwn-model-mabuchi}
\]
and
\[
 J^{\NA}(\cX,\mathcal A_D)
 :=\langle\mathcal A_D\rangle\cdot A^n
      -\frac1{n+1}\langle\mathcal A_D^{n+1}\rangle.
\tag{62.6}\label{eq:mpwn-model-j}
\]
All intersections are on \(\cX\).  The correction in
\eqref{eq:mpwn-model-mabuchi} replaces the central fiber by its reduction;
it is exactly the geometric manifestation of the reduced discrepancy
\(A_{X\times\mathbf P^1}-1\).

\begin{example}\label{ex:mpwn-big-translation-failure}
The expressions above need not be invariant under a base translation
of an arbitrary big class.  Let
\[
 \cX=\operatorname{Bl}_{(p,0)}
          (\mathbf P^1_x\times\mathbf P^1_t),
 \qquad A=H,
\]
where $H$ and $F$ are the pullbacks of the point classes from the first
and second factors, and $E$ is the exceptional curve.  Thus
\[
 H^2=F^2=H\cdot E=F\cdot E=0,\qquad
 H\cdot F=1,\qquad E^2=-1.
\]
The central fiber is $(F-E)+E$ and is reduced.  Set
$B_c=H+cF-\frac12E$.  For $0<c<\frac12$, its Zariski decomposition is
\[
 \begin{split}
 B_c&=P_c+N_c,\\
 P_c&=\left(\frac12+c\right)H+cF-cE,\\
 N_c&=\left(\frac12-c\right)(H-E).
 \end{split}
 \tag{62.6a}\label{eq:mpwn-big-translation-zariski}
\]
Indeed, the cone of effective curves of this toric surface is generated
by $E$, $H-E$, and $F-E$.  Their intersections with $P_c$ are
$c$, $0$, and $\frac12$, respectively, so $P_c$ is nef.
Moreover, $(H-E)^2=-1$ and $P_c\cdot N_c=0$.
Consequently
\[
 \langle B_c\rangle=P_c,\qquad
 \langle B_c^2\rangle=P_c^2=c+c^2>0.
\]
Here $V=1$ and $K_{\cX/\mathbf P^1}=-2H+E$.  Substitution into
\eqref{eq:mpwn-model-mabuchi} and \eqref{eq:mpwn-model-j} gives
\[
 M^{\NA}(\cX,B_c)=c^2,\qquad
 J^{\NA}(\cX,B_c)=\frac c2-\frac{c^2}{2}.
\]
At $c=\frac32$, the class $B_c$ is K\"ahler, so its positive products
are ordinary intersections.  The resulting values are
\begingroup
\renewcommand{\arraystretch}{1.8}
\[
 \begin{array}{c|cc}
 c&M^{\NA}(\cX,B_c)&J^{\NA}(\cX,B_c)\\ \hline
 \dfrac1{10}&\dfrac1{100}&\dfrac9{200}\\
 \dfrac32&\dfrac14&\dfrac18
 \end{array}
 \tag{62.6b}\label{eq:mpwn-big-translation-values}
\]
\endgroup
Writing $D_c=cF-\frac12E$, we have
$D_{3/2}=\frac32(F-E)+E$, so $f_{D_{3/2}}\geq1$.
The earlier divisor has a negative coefficient along $E$.
Although the two divisors differ by $\frac75F$, neither number is
unchanged.  In particular,
$F\cdot\langle B_{1/10}\rangle=\frac35\ne V$.
The vanishing $F^2=0$ does not license a binomial expansion of the
positive product of a non-nef big class.  The negative part in
\eqref{eq:mpwn-big-translation-zariski} is horizontal and varies with $c$.
\end{example}

\begin{definition}\label{def:mpwn-two-stabilities}
\begin{enumerate}
\item $(X,\alpha)$ is \emph{uniformly K-stable} if there is
      $\delta>0$ such that
      \[
          M^{\NA}(\cX,\mathcal A_D)\geq
          \delta J^{\NA}(\cX,\mathcal A_D)
      \tag{62.7}\label{eq:mpwn-ordinary-uniform-k}
      \]
      for every K\"ahler test configuration.
\item It is \emph{uniformly K-stable for models} if
      \eqref{eq:mpwn-ordinary-uniform-k} holds for every big test
      configuration.
\end{enumerate}
\end{definition}

These are the literal intersection definitions in
\cite[Definition~1.1.2]{MPWN25}.  The second condition is a priori
stronger because every K\"ahler class is big.  We retain this meaning
of stability for models and distinguish the following restricted test.

\begin{definition}\label{def:mpwn-normalized-model-test}
The \emph{normalized model inequality} holds if there is $\delta>0$
such that \eqref{eq:mpwn-ordinary-uniform-k} holds for every big
test configuration whose associated PL function satisfies $f_D\geq1$.
\end{definition}

Stability for all big models implies the normalized model inequality
with the same $\delta$.  The latter is the condition identified with
potential-theoretic coercivity below.

\section{From a normalized big model to a finite-energy potential}

Let $f_D\in\PL_{\mathbf R}(X^{\bethan})$ be represented by a vertical
real divisor $D$, and assume
\[
                              f_D\geq1.
\tag{62.8}\label{eq:mpwn-fd-positive}
\]
At each component of $\cX_0=\sum_i b_iE_i$, this says that the
coefficient of $D$ is at least $b_i$.  Thus $D-\cX_0$ is effective.
The class $A+[\cX_0]$ is the pullback of a K\"ahler class on
$X\times\mathbf P^1$, hence is big on $\cX$; adding the effective
class $D-\cX_0$ shows that $A+D$ is big as well.
The continuous envelope
\[
                              \phi_D:=P_\alpha(f_D)
\tag{62.9}\label{eq:mpwn-model-envelope}
\]
is bounded between the minimum and maximum of $f_D$, so it belongs
to $\mathcal E^1(\alpha)$.

\begin{theorem}\label{thm:mpwn-model-intersection-formulas}
For a big test configuration satisfying
\eqref{eq:mpwn-fd-positive},
\[
 M_\alpha^{\NA}(P_\alpha(f_D))
      =V^{-1}M^{\NA}(\cX,A+D),
\qquad
 J_\alpha(P_\alpha(f_D))
      =V^{-1}J^{\NA}(\cX,A+D).
\tag{62.10}\label{eq:mpwn-model-functional-identities}
\]
\end{theorem}

\begin{proof}
We use the restricted-positive-product identities of
\cite[Propositions~9.1.2, 9.1.3, and 9.1.5]{MPWN25}
in their range $f_D\geq1$.  They give
\[
 \MA_\alpha(P_\alpha(f_D))
 =V^{-1}\sum_i b_i
       \langle(A+D)^n\rangle_{\cX\mid E_i}\delta_{v_{E_i}}
\tag{62.11}\label{eq:mpwn-big-envelope-ma}
\]
and
\[
 \langle(A+D)^n\rangle_{\cX\mid E_i}
       =\langle(A+D)^n\rangle\cdot E_i.
\tag{62.12}\label{eq:mpwn-restricted-positive}
\]
Integration of reduced discrepancy therefore gives the entropy term
\[
 V^{-1}\langle(A+D)^n\rangle\cdot
 K^{\log}_{\cX/(X\times\mathbf P^1)},
\]
where the log correction is $\cX_{0,\mathrm{red}}-\cX_0$.
The same imported identities yield
\[
 E_\alpha(P_\alpha(f_D))
    =\frac{V^{-1}}{n+1}\langle(A+D)^{n+1}\rangle
\tag{62.13}\label{eq:mpwn-envelope-energy-positive-product}
\]
and the twisted-energy identity.  Their proof uses movable
approximants on modifications.  The nonnegative model potentials
in that construction have effective vertical divisors; relative
nefness then gives nefness on the total space by
\cite[Lemma~9.1.4]{MPWN25}.
One does not approximate a non-nef big class by K\"ahler classes
having its same minimal singularities.

Combining energy, twisted energy, and reduced entropy gives the
Mabuchi identity; the imported mixed formula with one positive
factor gives the $J$ identity.  These are
\cite[Theorem~9.1.1]{MPWN25}, with the coefficient correction at
\eqref{eq:mpwn-s-constant}.
\end{proof}

\begin{proposition}\label{prop:mpwn-invariant-invariance}
Suppose $f_D\geq1$.  The quantities
$M^{\NA}(\cX,A+D)$ and $J^{\NA}(\cX,A+D)$ are unchanged under
pullback to a dominating smooth model.
They are also unchanged by $D\mapsto D+c\cX_0$ whenever
both $f_D$ and $f_D+c$ are at least one.
\end{proposition}

\begin{proof}
Pulling back a vertical divisor leaves its PL function unchanged.
Adding $c\cX_0$ replaces that function by $f_D+c$, and the envelope
satisfies $P_\alpha(f_D+c)=P_\alpha(f_D)+c$.
Apply \Cref{thm:mpwn-model-intersection-formulas} to each representative.
The potential-theoretic functionals $M_\alpha^{\NA}$ and $J_\alpha$
are invariant under constants: the changes in the energy and
twisted-energy terms cancel by \eqref{eq:mpwn-s-constant},
while entropy depends only on the unchanged Monge--Amp\`ere measure.
This proves both assertions without a pullback or binomial rule
for arbitrary positive products.
\end{proof}

\begin{corollary}\label{cor:mpwn-model-stability-pl-envelopes}
The normalized model inequality of
\Cref{def:mpwn-normalized-model-test} is equivalent to the existence of
$\delta>0$ such that
\[
 M_\alpha^{\NA}(P_\alpha(f))
       \geq\delta J_\alpha(P_\alpha(f))
\tag{62.14}\label{eq:mpwn-pl-envelope-stability}
\]
for every $f\in\PL_{\mathbf R}(X^{\bethan})$.
Uniform K-stability for all big models implies these conditions.
\end{corollary}

\begin{proof}
Represent $f$ by a vertical real divisor $D$ and choose
$c\geq1-\min f$.  Then $f+c\geq1$, and
$A+D+c[\cX_0]$ is big by the effectiveness argument above.
Apply the normalized model inequality to this class and then use
\eqref{eq:mpwn-model-functional-identities}.
Translation invariance of the potential functionals gives
\eqref{eq:mpwn-pl-envelope-stability} for the original $f$.
Conversely, apply the PL-envelope inequality to every $f_D\geq1$
and use the same identities.  Both directions preserve $\delta$.
The final implication follows by restricting the set of test
configurations in \Cref{def:mpwn-two-stabilities}.
\end{proof}

\begin{remark}\label{rem:mpwn-unnormalized-model-boundary}
\cite[Corollary~9.1.6]{MPWN25} asserts the equivalence with the
all-big-model definition.  The preceding argument proves the
equivalence for the normalized model inequality and the implication
from the all-big-model condition.  Its converse for the literal
raw expressions \eqref{eq:mpwn-model-mabuchi} and
\eqref{eq:mpwn-model-j} requires an additional comparison.
It does not follow by adding a base class:
\Cref{ex:mpwn-big-translation-failure} disproves that invariance.
The example refutes this proposed proof step, not by itself the
equivalence of the two stability conditions.  No such converse is
used below.
\end{remark}

\section{Completion from models to all finite energy}

Call uniform stability on the full finite-energy space
\[
 M_\alpha^{\NA}(\phi)\geq\delta J_\alpha(\phi),
 \qquad\phi\in\mathcal E^1(\alpha),
\tag{62.15}\label{eq:mpwn-bk-stability}
\]
uniform \(\widehat K\)-stability (also called uniform
\(b\)-K-stability in \cite{MP24,MPWN25}).  This is not a definition by
test configurations; it is coercivity of a functional on the completed
potential space.

\begin{theorem}\label{thm:mpwn-model-e1-equivalence}
The following conditions are equivalent:
\begin{enumerate}
\item there is \(\delta>0\) such that
      \eqref{eq:mpwn-bk-stability} holds on all
      \(\mathcal E^1(\alpha)\);
\item it holds for every continuous \(\alpha\)-psh \(\phi\) whose
      Monge--Amp\`ere measure is supported on one dual complex;
\item it holds for every real PL envelope \(P_\alpha(f)\), with
      \(f\in\PL_{\mathbf R}\).
\end{enumerate}
Consequently, the normalized model inequality is equivalent to uniform
$\widehat K$-stability.  Uniform K-stability for all big models implies
these equivalent conditions.
\end{theorem}

\begin{proof}[Imported entropy-approximation core]
The first condition trivially implies the second.  For the converse, take
\(\phi\in\mathcal E^1(\alpha)\) and put \(\mu=\MA_\alpha(\phi)\).
By \Cref{thm:mpwn-measure-approximation}, choose measures
\(\mu_j\), supported on dual complexes, that converge strongly to \(\mu\)
and satisfy \(\Ent(\mu_j)\to\Ent(\mu)\).  Let
\(\MA_\alpha(\phi_j)=\mu_j\), with sup normalization.  The Calabi--Yau
homeomorphism gives \(\phi_j\to\phi\) strongly, dual-complex regularity gives
\(\phi_j\in C^0\), and entropy convergence yields
\[
 J_\alpha(\phi_j)\to J_\alpha(\phi),
 \qquad
 M_\alpha^{\NA}(\phi_j)\to M_\alpha^{\NA}(\phi).
\]
Pass the assumed inequality to the limit.

It remains to pass from PL envelopes to continuous potentials supported on
one complex.  If
\(\operatorname{supp}\MA_\alpha(\phi)\subset\Delta_{\cX}\), locality gives
\(\phi=P_\alpha(\phi\circ p_{\cX})\).  Approximate the continuous convex
function \(\phi|_{\Delta_{\cX}}\) uniformly by rational PL functions on
subdivisions.  Their envelopes converge uniformly, hence strongly, to
\(\phi\); all their measures remain on the same underlying monomial
complex.  Weak convergence gives convergence of every fixed
dual-complex discrepancy integral.  Since all these measures are supported
on the same monomial complex, reduced discrepancy there equals its affine
model representative; hence these fixed integrals are their entropies.  The
functional inequality passes to the limit.

The equivalence of the third condition with the normalized model
inequality is \Cref{cor:mpwn-model-stability-pl-envelopes}.
The simultaneous strong and entropy approximation and the support
formula for envelopes are the imported cores from Sections~7--8 of
\cite{MPWN25}; they give the potential-theoretic equivalence in
\cite[Theorem~9.2.1]{MPWN25}.
The implication from stability for all big models follows by
restriction to those satisfying $f_D\geq1$.
\end{proof}

\begin{corollary}\label{cor:mpwn-model-stability-csck}
If \((X,\alpha)\) is uniformly K-stable for models, then it satisfies the
finite-energy coercivity condition \eqref{eq:mpwn-bk-stability}; hence the
later-source cscK theorem produces a unique cscK metric in \(\alpha\).
\end{corollary}

\begin{proof}[Imported conclusion and formal deduction]
Restrict the all-big-model hypothesis to $f_D\geq1$ and apply
\Cref{thm:mpwn-model-e1-equivalence} to obtain finite-energy coercivity.
This implication needs no comparison with the raw invariants of an
unnormalized big class.  The conclusion of
\cite[Theorem~2.4.3]{MPWN25} then yields a unique cscK metric in \(\alpha\).
Equivalently, the entire implication is stated directly as
\cite[Corollary~9.2.2]{MPWN25}.  Both cited statements use the reduced
entropy of \cite[Definition~2.4.1]{MPWN25}.
The reduced entropy bound is also available in
\Cref{thm:mp-reduced-entropy-slope}; the present cscK conclusion uses the
cited existence theorem directly.
\end{proof}

\begin{remark}[Scope of the cited conclusion]
\cite[Theorem~2.4.3 and Corollary~9.2.2]{MPWN25} state literal uniqueness
under the ordinary \(J\)-coercive hypotheses used here.
Automorphism-relative stability is a different statement and is not
inferred from these passages.
\end{remark}

\section{The beta invariant on measures}

The K\"ahler cone is open, so \(\alpha+tK_X\) is K\"ahler for all
sufficiently small \(|t|\) and \(K_X\) is an admissible tangent direction at
\(\alpha\).  It is not itself a K\"ahler class, and nothing below requires
it to be.  For \(\mu\in\mathcal M^1(\alpha)\), define
\[
 \beta_\alpha(\mu)
 :=\Ent(\mu)+\nabla_{K_X}E_\alpha^\vee(\mu)
 :=\Ent(\mu)+
    \left.\frac{\rd}{\rd t}\right|_{t=0}
       E_{\alpha+tK_X}^\vee(\mu).
\tag{62.16}\label{eq:mpwn-beta}
\]
Differentiability in the class variable is not proved here.  It is the
single imported analytic input of this section: every statement from
\eqref{eq:mpwn-beta} onward, including
\Cref{thm:mpwn-beta-mabuchi} and \Cref{thm:mpwn-valuative-criterion},
presupposes that \(t\mapsto E_{\alpha+tK_X}^\vee(\mu)\) is differentiable at
\(t=0\).  It is an imported theorem from the
finite-energy measure formalism.  The entropy uses
\(A_{X\times\mathbf P^1}-1\), not total discrepancy.

The definition is the transcendental counterpart of the divisorial-measure
formalism of \cite{BJKS2}.  That source works with polarized projective
varieties and algebraic norms.  Here the class derivative is taken in the
K\"ahler cone and the measure space comes from analytic models; no projective
specialization is used in the proof below.

\begin{theorem}\label{thm:mpwn-beta-mabuchi}
For every \(\phi\in\mathcal E^1(\alpha)\),
\[
 \beta_\alpha(\MA_\alpha(\phi))
      =M_\alpha^{\NA}(\phi),
\qquad
 E_\alpha^\vee(\MA_\alpha(\phi))
      \asymp J_\alpha(\phi),
\tag{62.17}\label{eq:mpwn-beta-mabuchi}
\]
where the comparison constants depend only on \(n\).
\end{theorem}

\begin{proof}[Imported class-derivative core]
For a model potential, write the dual energy as an intersection number and
differentiate it with respect to the background class.  The derivative equals
\(sE_\alpha+E_\alpha^{K_X}\); adding entropy gives Mabuchi.  Approximate an
arbitrary finite-energy potential decreasingly by model potentials and use
continuity of the class derivative, as in
\cite[Proposition~4.5]{BJ25FiniteEnergy}.  The second assertion follows from
\[
 E_\alpha^\vee(\MA_\alpha(\phi))
 =E_\alpha(\phi)-\int\phi\,\MA_\alpha(\phi)
 =J_\alpha(\phi,0)
\]
and quasi-symmetry of the Dirichlet energy.  This is
\cite[Theorem~10.0.1]{MPWN25}.
\end{proof}

\begin{definition}[Divisorial measures]
\label{def:mpwn-divisorial-measures}
Let \(\mathcal M^{\divis}\) consist of probability measures
\[
                       \mu=\sum_{i=1}^N a_i\delta_{v_i},
 \qquad a_i\geq0,\quad\sum_i a_i=1,\quad v_i\in X^{\divis}.
\tag{62.18}\label{eq:mpwn-divisorial-measure}
\]
\end{definition}

\begin{theorem}\label{thm:mpwn-valuative-criterion}
The normalized model inequality of
\Cref{def:mpwn-normalized-model-test} is equivalent to the existence of
$\delta>0$ such that
\[
             \beta_\alpha(\mu)
             \geq\delta E_\alpha^\vee(\mu)
        \quad\text{for every }\mu\in\mathcal M^{\divis}.
\tag{62.19}\label{eq:mpwn-beta-criterion}
\]
Equivalently,
\[
 \inf_{\substack{\mu\in\mathcal M^{\divis}\\
                  E_\alpha^\vee(\mu)>0}}
       \frac{\beta_\alpha(\mu)}{E_\alpha^\vee(\mu)}>0.
\tag{62.20}\label{eq:mpwn-beta-ratio}
\]
Uniform K-stability for all big models implies these conditions.
\end{theorem}

\begin{proof}
The Calabi--Yau theorem represents every \(\mu\in\mathcal M^1\) as
\(\MA_\alpha(\phi)\); \Cref{thm:mpwn-beta-mabuchi} translates
finite-energy stability into the beta inequality for all finite-energy
measures.  For a divisorial measure, the solution is a PL envelope by
\Cref{cor:mpwn-divisorial-envelope}.  Conversely, the
Monge--Amp\`ere measure of a PL envelope is divisorial.  Thus the PL-envelope
test of \Cref{thm:mpwn-model-e1-equivalence} is equivalent to
\eqref{eq:mpwn-beta-criterion}.  The equivalence is not an equality of
thresholds: \Cref{thm:mpwn-beta-mabuchi} gives
\(\beta_\alpha(\MA_\alpha\phi)=M^{\NA}_\alpha(\phi)\) exactly, but only
\(E^\vee_\alpha(\MA_\alpha\phi)\asymp J_\alpha(\phi)\), so the constant
\(\delta\) in \eqref{eq:mpwn-beta-criterion} differs from the one in the
normalized model test by the $n$-dependent comparison constants.
This proves the valuative equivalence in
\cite[Theorem~10.0.4]{MPWN25} with the normalized model condition.
The distinction from the raw all-big-model condition is
\Cref{rem:mpwn-unnormalized-model-boundary}.
\end{proof}

\begin{warning}{Measures, not only single divisors}
The criterion ranges over finite weighted combinations of divisorial
valuations.  Positivity only for Dirac masses is a different, generally
weaker test.  Dreaminess is not assumed in
\eqref{eq:mpwn-beta-criterion}.
\end{warning}

\section{Finite-dimensional volume formula}

Let \(F_1,\ldots,F_\ell\) be prime divisors on one modification
\(\rho:Y\to X\), let \(v_i=\ord_{F_i}\), and take a rational probability
vector \(\xi=(\xi_i)\).  Put
\[
                         \mu_\xi:=\sum_{i=1}^\ell\xi_i\delta_{v_i}.
\tag{62.21}\label{eq:mpwn-muxi}
\]
For \(t=(t_i)\in\mathbf R^\ell\), define the constrained envelope
\[
 \phi_t
 :=\sup\{\phi\in\PSH(\alpha)\mid
                    \phi(v_i)\leq t_i,\ 1\leq i\leq\ell\},
\qquad
 f_\alpha(t):=E_\alpha(\phi_t).
\tag{62.22}\label{eq:mpwn-constrained-envelope}
\]
The supremum means its psh envelope; by the finite-constraint lemma it is
continuous.  Convex duality gives
\[
 E_\alpha^\vee(\mu_\xi)
    =\sup_{t\in\mathbf R^\ell}
       \{f_\alpha(t)-\xi\cdot t\}.
\tag{62.23}\label{eq:mpwn-energy-legendre}
\]

The probability constraint is visible in the translation direction.  If
\(\mathbf1=(1,\ldots,1)\), then the candidate sets in
\eqref{eq:mpwn-constrained-envelope} give
\[
            \phi_{t+c\mathbf1}=\phi_t+c,
 \qquad     f_\alpha(t+c\mathbf1)=f_\alpha(t)+c.
\]
Since \(\xi\cdot\mathbf1=1\), the expression
\(f_\alpha(t)-\xi\cdot t\) is unchanged by the same translation.  Thus the
Legendre problem has exactly one redundant additive direction, matching the
normalization of potentials modulo constants.

\begin{theorem}\label{thm:mpwn-volume-formula}
For any \(\lambda_0<\min_i t_i\),
\[
 f_\alpha(t)
 =\lambda_0+\frac1V\int_{\lambda_0}^{+\infty}
     \left\langle
       \left(\rho^*\alpha-
         \sum_{i=1}^{\ell}(\lambda-t_i)_+F_i\right)^n
     \right\rangle\,\rd \lambda.
\tag{62.24}\label{eq:mpwn-volume-formula}
\]
The positive product is declared zero after the class ceases to be
pseudoeffective, so the improper integral is finite.
\end{theorem}

\begin{proof}[Imported test-curve energy formula]
The family of maximal \(\omega\)-psh functions
\[
 G_\lambda
 :=\sup\{u\in\PSH_{\sup}(X,\omega)\mid
              \nu_{F_i}(u)\geq\lambda-t_i\ \forall i\}
\tag{62.25}\label{eq:mpwn-rwn-test-curve}
\]
is a concave \(\mathcal I\)-maximal test curve.  Its Ross--Witt Nystr\"om
transform is the maximal ray associated to \(\phi_t\)
\cite[Lemma~10.1.1]{MPWN25}.  The DXZ energy formula
\cite[Theorem~2.6(iii)]{DXZ23} computes the ray slope by integrating the
non-pluripolar masses of the \(G_\lambda\).
Removing their divisorial singularities identifies that mass with
\[
 \left\langle
   \left(\rho^*\alpha-
       \sum_i(\lambda-t_i)_+F_i\right)^n
 \right\rangle.
\]
Let \(\lambda_{\max}\) be the terminal parameter of the test curve and let
\(m(\lambda)\) denote the displayed positive product.  The precise DXZ
identity is
\[
 f_\alpha(t)=\lambda_{\max}
   +\int_{-\infty}^{\lambda_{\max}}
        \bigl(-1+V^{-1}m(\lambda)\bigr)\,\rd\lambda.
\]
For \(\lambda\leq\lambda_0<\min_i t_i\), all positive parts vanish and
\(m(\lambda)=V\), so the integrand is zero.  Therefore
\[
 \begin{split}
 f_\alpha(t)
 &=\lambda_{\max}
   +\int_{\lambda_0}^{\lambda_{\max}}
       \bigl(-1+V^{-1}m(\lambda)\bigr)\,\rd\lambda\\
 &=\lambda_0+V^{-1}
       \int_{\lambda_0}^{\lambda_{\max}}m(\lambda)\,\rd\lambda.
 \end{split}
\]
After \(\lambda_{\max}\) the class is not pseudoeffective and its positive
product is zero, so extending the upper limit to \(+\infty\) gives
\eqref{eq:mpwn-volume-formula}.  This is
\cite[Theorem~10.1.2]{MPWN25}.  This use of \cite{DXZ23} is on the proved
\(\mathcal I\)-maximal finite-energy locus only.
\end{proof}

If \(g_\alpha(\xi)\) denotes the Legendre transform in
\eqref{eq:mpwn-energy-legendre}, then
\[
 \beta_\alpha(\mu_\xi)
 =\sum_{i=1}^{\ell}\xi_iA_X(F_i)
      +\nabla_{K_X}g_\alpha(\xi).
\tag{62.26}\label{eq:mpwn-beta-explicit}
\]
This is \cite[Corollary~10.1.4]{MPWN25}.

\begin{example}[One divisor]
\label{ex:mpwn-one-divisor-beta}
For \(\mu=\delta_{\ord_F}\), translation invariance reduces
\eqref{eq:mpwn-volume-formula} to
\[
 E_\alpha^\vee(\delta_{\ord_F})
   =\frac1V\int_0^{+\infty}
          \langle(\rho^*\alpha-\lambda F)^n\rangle\,\rd \lambda.
\tag{62.27}\label{eq:mpwn-one-divisor-energy}
\]
Indeed, normalize the one-point constrained envelope by
\(\phi_0(\ord_F)=0\).  Orthogonality makes it the variational maximizer for
\(\delta_{\ord_F}\), so its energy equals the dual energy, and
\eqref{eq:mpwn-volume-formula} with \(t=0\) gives the displayed integral.
The integrand vanishes beyond the pseudoeffective threshold.  Formula
\eqref{eq:mpwn-beta-explicit} then differentiates this integral in the
\(K_X\)-direction and adds \(A_X(F)\).  The signs match
\(\rho^*\alpha-\lambda F\): imposing more vanishing consumes positivity.
\end{example}

\begin{example}[Toric specialization]
\label{ex:mpwn-toric-beta}
If \(X\) is a smooth projective toric manifold and the \(F_i\) are toric
prime divisors on a toric modification, the positive intersection in
\eqref{eq:mpwn-volume-formula} is \(n!\) times the Euclidean volume of the
moment polytope cut by the inequalities corresponding to
\((\lambda-t_i)_+\).  Thus \(f_\alpha\) is an integral of sliced polytope
volumes, and \eqref{eq:mpwn-energy-legendre} becomes ordinary finite
dimensional convex duality.  This recovers the toric sign convention of
\Cref{chap:transcendental-okounkov-dh} without asserting that every analytic model is toric.
\end{example}

\section{Flag ideals, envelopes, and radial rays}

For a model represented by a flag ideal and vertical divisor $D$
with $f_D\geq1$, the proved chain of constructions is
\[
 \begin{tikzcd}[column sep=large]
  (\cX,A+D)
 \arrow[r,mapsto,"f_D"]
 & f_D\in\PL_{\mathbf R}
   \arrow[r,mapsto,"P_\alpha"]
 & \phi_D\in\mathcal E^1(\alpha)
   \arrow[r,mapsto,"\text{maximal ray}"]
 & U^{\phi_D}.
 \end{tikzcd}
\tag{62.28}\label{eq:mpwn-model-envelope-ray-chain}
\]
The first arrow retains the vertical ideal data; the second discards the
non-nef part by an obstacle problem; the last uses the radial
correspondence.  The proved numerical bridge to the associated ray is
\[
 J^{\NA}(\cX,A+D)
   =V\,J_\alpha(\phi_D)
   =V\lim_{s\to\infty}\frac{J_\omega(U^{\phi_D}_s)}s.
\tag{62.29}\label{eq:mpwn-model-ray-slopes}
\]
The model identity
$M^{\NA}(\cX,A+D)=V M_\alpha^{\NA}(\phi_D)$ is
\eqref{eq:mpwn-model-functional-identities}.  Together with
\Cref{thm:mp-reduced-entropy-slope}, it gives the companion inequality
\[
 M^{\NA}(\cX,A+D)
 \leq V\liminf_{s\to\infty}\frac{M_\omega(U^{\phi_D}_s)}s.
 \tag{62.29a}\label{eq:mpwn-model-mabuchi-slope}
\]
This is a lower bound for the radial slope.  Equality would require a
separate entropy-slope equality and is not inferred here.
\Cref{cor:mpwn-model-stability-csck} still uses the later primary-source
existence theorem directly.

\section{The theorem boundary as of September 2026}

The following distinctions summarize what is used in this book.

\begingroup
\setlength{\parindent}{0pt}

\textbf{Statement}\par\nopagebreak
Finite-energy uniform stability
\(\Rightarrow\) cscK\par\nopagebreak

\textbf{Status used here}\par\nopagebreak
Imported from \cite[Theorem~2.4.3]{MPWN25}, which explicitly uses reduced
entropy.  The reduced entropy lower bound is proved in
\Cref{thm:mp-reduced-entropy-slope}; it is distinct from the deep cscK
existence implication used here.\par\nopagebreak

\textbf{What it does not imply}\par\nopagebreak
No converse; automorphism-relative variants require extra formulation.

\medskip
\textbf{Statement}\par\nopagebreak
Normalized model inequality
$\Leftrightarrow$ finite-energy uniform stability\par\nopagebreak

\textbf{Status used here}\par\nopagebreak
\Cref{thm:mpwn-model-e1-equivalence}, using
\cite[Theorem~9.2.1]{MPWN25}.\par\nopagebreak

\textbf{What it does not imply}\par\nopagebreak
The converse to stability for every raw big class requires the
additional comparison recorded in
\Cref{rem:mpwn-unnormalized-model-boundary}.

\medskip
\textbf{Statement}\par\nopagebreak
Model uniform stability
\(\Rightarrow\) cscK\par\nopagebreak

\textbf{Status used here}\par\nopagebreak
Imported from \cite[Corollary~9.2.2]{MPWN25}; its formal reduction uses
\cite[Theorems~9.2.1 and 2.4.3]{MPWN25}.\par\nopagebreak

\textbf{What it does not imply}\par\nopagebreak
Not an unrestricted ordinary compact-K\"ahler YTD equivalence.

\medskip
\textbf{Statement}\par\nopagebreak
Ordinary uniform stability using only K\"ahler test configurations\par\nopagebreak

\textbf{Status used here}\par\nopagebreak
A distinct, a priori weaker condition in
\cite[Definition~1.1.2]{MPWN25}.\par\nopagebreak

\textbf{What it does not imply}\par\nopagebreak
The model-completion argument cannot be applied without control of big
configurations.

\endgroup

Two 2026 preprints lie beyond the dependency-checked core above.
Boucksom--Jonsson's version~2 overview, dated 19 May 2026,
\cite{BJ26ComplexNA}, describes a broader complex-to-non-Archimedean YTD
program; we use it for orientation, not as a proof source.  Trusiani's
version~2 preprint, dated 17 August 2026, \cite{Trusiani26}, retains its
reported full
automorphism-relative uniform YTD theorem for smooth polarized
\emph{projective} varieties together with entropy regularization, and adds
an appendix by S\'{e}bastien Boucksom.  In the present book draft these are
recorded as dated research claims: their hypotheses and dependencies have
not been folded into the proofs above, and even if confirmed they would not
by themselves establish the corresponding theorem for an arbitrary
nonprojective compact K\"ahler class.

\section*{Source note}

Definitions and positive-intersection invariants are MPWN25
Sections~1.1 and 2.1.  The exact envelope/intersection identities are
Section~9.1, with the independently verified coefficient correction isolated
at \eqref{eq:mpwn-s-constant}.  Finite-energy stability implies cscK is
\cite[Theorem~2.4.3]{MPWN25}; equivalence of model and finite-energy
stability and the model cscK implication are Section~9.2.  These conclusions
are imported from that later source.  The reduced entropy lower bound
proved in \Cref{thm:mp-reduced-entropy-slope} also gives
\eqref{eq:mpwn-model-mabuchi-slope}, but does not replace the imported
cscK existence theorem.  The measure-theoretic beta invariant and valuative
criterion are Section~10.0; the constrained-envelope volume formula is
Section~10.1, whose ray-energy input is
\cite[Theorem~2.6(iii)]{DXZ23}.  In the polarized projective setting, the
measure-valuative architecture should be compared with \cite{BJKS2}; the
present theorem remains the analytic compact-K\"ahler statement of
\cite{MPWN25}.  The final status map is deliberately dated August~2026 and
separates dependency-audited theorems from recent announcements.
